\section{هفته‌ی 5: عدم قطعیت برآوردها و برازش برای چند شرکت‌کننده}

\begin{quote}
\textbf{فصل 2: روش‌های برآورد پارامتر} · درس 1 ساعت · \lr{lab} 2 ساعت
\textbf{داده‌ی اصلی:} داده‌های فراموشی هفته‌های 2\lr{--}4 (برآوردهای ما چقدر
قطعی‌اند؟) و یک آزمایش \textbf{تمرین} \lr{(practice)} شبیه‌سازی‌شده‌ی جدید
با 20 شرکت‌کننده (قانون توانی تمرین).
\end{quote}

\begin{center}\rule{0.5\linewidth}{0.5pt}\end{center}

\subsection{واژه‌های کلیدی}

واژه‌های مهم این درس با معادل فارسی‌شان (حتی اگر در درس‌های قبل آمده
باشند).

{\def\LTcaptype{none} % do not increment counter
\begin{longtable}[]{@{}
  >{\raggedright\arraybackslash}p{(\linewidth - 4\tabcolsep) * \real{0.3333}}
  >{\raggedright\arraybackslash}p{(\linewidth - 4\tabcolsep) * \real{0.3333}}
  >{\raggedright\arraybackslash}p{(\linewidth - 4\tabcolsep) * \real{0.3333}}@{}}
\toprule\noalign{}
\begin{minipage}[b]{\linewidth}\raggedright
\lr{English}
\end{minipage} & \begin{minipage}[b]{\linewidth}\raggedright
فارسی
\end{minipage} & \begin{minipage}[b]{\linewidth}\raggedright
توضیح کوتاه
\end{minipage} \\
\midrule\noalign{}
\endhead
\bottomrule\noalign{}
\endlastfoot
\lr{estimate} / \lr{estimator} & برآورد / برآوردگر & مقدار به‌دست‌آمده
برای پارامتر / روش به دست آوردن آن \\
\lr{fit} / \lr{fitting} & برازش & تنظیم پارامترها تا مدل با داده جور
شود \\
\lr{parameter} & پارامتر & کمیت قابل تنظیم مدل که بین افراد یا شرایط فرق
می‌کند \\
\lr{simulation} & شبیه‌سازی & تولید داده‌ی مصنوعی با اجرای مدل \\
\lr{power law of practice} & قانون توانی تمرین & کاهش توانی زمان واکنش
با تمرین \\
\lr{exponential} / \lr{power function} & تابع نمایی / توانی & دو شکل
رایج برای منحنی فراموشی و یادگیری \\
\lr{sampling distribution} & توزیع نمونه‌گیری & توزیع برآورد در تکرارهای
فرضی آزمایش \\
\lr{standard error (SE)} & خطای استاندارد & پراکندگی برآورد در تکرارهای
فرضی آزمایش \\
\lr{confidence interval (CI)} & بازه‌ی اطمینان & بازه‌ای که با روش مشخص
پارامتر واقعی را در بر می‌گیرد \\
\lr{Hessian} & ماتریس \lr{Hessian} & ماتریس مشتق‌های دوم؛ انحنای
درست‌نمایی \\
\lr{log-likelihood} / \lr{NLL} & لگاریتم درست‌نمایی / منفی آن & مقداری که
بهینه‌سازها کمینه یا بیشینه می‌کنند \\
\lr{Wald interval} & بازه‌ی \lr{Wald} & برآورد \ensuremath{\pm} 1.96 \lr{SE} \\
\lr{profile likelihood} & درست‌نمایی نیم‌رخ \lr{(profile)} & درست‌نمایی یک
پارامتر با بهینه کردن بقیه \\
\lr{likelihood} & درست‌نمایی & احتمال داده به‌عنوان تابعی از پارامترها \\
\lr{bootstrap} & \lr{bootstrap} (بازنمونه‌گیری) & تخمین عدم قطعیت با
شبیه‌سازی یا بازنمونه‌گیری \\
\lr{parameter trade-off} & بده‌بستان پارامترها & تغییر یک پارامتر که با
پارامتر دیگر جبران می‌شود \\
\lr{aggregation artifact} & خطای ناشی از میانگین‌گیری & الگویی در میانگین
که در هیچ فردی نیست \\
\lr{trial} & کوشش \lr{(trial)} & یک بار انجام تکلیف توسط شرکت‌کننده \\
\lr{bounds} / \lr{constraints} & کران‌ها / قیدها & محدودیت‌های مجاز برای
مقدار پارامترها \\
\lr{accuracy} & دقت & نسبت پاسخ‌های درست \\
\lr{complete} / \lr{no} / \lr{partial pooling} & ادغام کامل / بدون ادغام
/ ادغام جزئی \lr{(pooling)} & یک پارامتر برای همه / جدا برای هر فرد /
سلسله‌مراتبی \\
\lr{individual differences} & تفاوت‌های فردی & فرق پارامترها بین افراد \\
\lr{noise} & نویز & تغییرات تصادفی در شواهد یا در داده \\
\lr{benchmark} & \lr{benchmark} (مجموعه‌آزمون استاندارد) & مجموعه سؤال
برای مقایسه‌ی مدل‌های هوش مصنوعی \\
\end{longtable}
}

\begin{center}\rule{0.5\linewidth}{0.5pt}\end{center}

\subsection{اهداف یادگیری}

در پایان این هفته، دانشجویان می‌توانند:

\begin{enumerate}
\def\labelenumi{\arabic{enumi}.}
\tightlist
\item
  توزیع نمونه‌گیری \lr{(sampling distribution)} یک برآوردگر را
  \textbf{توضیح دهند} و خطای استاندارد \lr{(SE)} را از انحراف معیار
  داده‌ها \textbf{متمایز کنند}.
\item
  خطاهای استاندارد و بازه‌های اطمینان پارامترهای مدل را به سه روش
  \textbf{حساب کنند}: از \lr{Hessian} تابع \lr{NLL} (روش \lr{Wald)}، از
  \lr{profile likelihood} و با \lr{bootstrap} پارامتری؛ و این سه روش را
  \textbf{مقایسه کنند}.
\item
  \lr{trade-off} پارامترها را از همبستگی برآوردشده بین پارامترها
  \textbf{تشخیص دهند}.
\item
  با شبیه‌سازی \textbf{نشان دهند} که میانگین‌گیری از منحنی‌های فردی می‌تواند
  شکل تابعی‌ای (یک «قانون توانی») بسازد که هیچ فردی از آن پیروی نمی‌کند.
\item
  بین برازش تجمیعی، فردی و (پیش‌نمایش) سلسله‌مراتبی \textbf{انتخاب کنند} و
  انتخابشان را \textbf{توجیه کنند}.
\end{enumerate}

\begin{center}\rule{0.5\linewidth}{0.5pt}\end{center}

\subsection{مقدمه}

\textbf{«\(\hat d = 0.42\)» نصف یک نتیجه است.} هفته‌ی قبل \lr{MLE} مدل
توانی برای داده‌های هفته‌ی 2 برابر \(\hat d = 0.42\) بود (مقدار واقعی
0.5). آیا 0.42 با 0.5 سازگار است؟ با 0.3؟ با 0.6؟ بدون یک معیار عدم
قطعیت نمی‌توانیم بگوییم دو شرکت‌کننده، دو گروه یا دو شرط واقعاً با هم فرق
دارند یا نه. از این به بعد هر برآورد پارامتری که گزارش می‌کنیم با یک
بازه‌ی اطمینان همراه است.

\textbf{قانونی که هیچ‌کس از آن پیروی نمی‌کند.} \emph{قانون توانی تمرین}
\lr{(power law of practice)} یکی از پرارجاع‌ترین یافته‌ها در روان‌شناسی
شناختی است: زمان واکنش با تمرین تقریباً به شکل تابع توانی از تعداد
\lr{trial}ها کم می‌شود. بعضی پژوهشگران استدلال کرده‌اند که شکل توانی
می‌تواند از \textbf{میانگین‌گیری} روی افرادی (یا آیتم‌هایی) بیاید که هرکدام
به‌طور \emph{نمایی} ولی با نرخ‌های متفاوت بهتر می‌شوند (مطالعه‌ی بیشتر را
ببینید). اگر این‌طور باشد، «قانون» میانگین را توصیف می‌کند ولی هیچ
یادگیرنده‌ی تکی را نه، و نظریه‌ای که برای توضیح یک تابع توانی در ذهن هر
فرد ساخته شود در واقع یک خطای ناشی از میانگین‌گیری را توضیح می‌دهد. این
هفته این اثر را با داده‌ی شبیه‌سازی‌شده نشان می‌دهیم و یاد می‌گیریم کی باید
افراد را برازش کرد نه میانگین گروه را.

\paragraph{این موضوع کجا به کار می‌رود؟}

\begin{itemize}
\tightlist
\item
  \textbf{هر ادعای تفاوت} («بیماران سریع‌تر از گروه کنترل فراموش می‌کنند»،
  «آموزش حساسیت را بالا برد») به برآورد عدم قطعیت نیاز دارد. بدون آن،
  تفاوت 0.05 در یک پارامتر ممکن است فقط نویز باشد.
\item
  \textbf{ارزیابی بالینی یک فرد:} برای اینکه بگوییم سرعت فراموشی یک
  بیمار غیرعادی است، به خطای استاندارد برآورد \emph{همان بیمار} نیاز
  داریم، نه فقط میانگین گروه.
\item
  \textbf{پژوهش یادگیری و مهارت:} «قانون توانی تمرین» مدت‌ها یک قانون
  همگانی حساب می‌شد. خطای ناشی از میانگین‌گیری این هفته نشان می‌دهد چرا
  امروز پژوهشگرها یادگیرنده‌ها را جداگانه برازش می‌کنند (و چرا همین احتیاط
  درباره‌ی منحنی‌های میانگین گروه در آموزش، علوم ورزشی و ارگونومی هم لازم
  است).
\item
  \textbf{تکرار پژوهش و فراتحلیل \lr{(meta-analysis)}} برآوردهای چند
  مطالعه را با وزن دادن به هر کدام بر اساس دقتش \lr{(1/SE²)} ترکیب
  می‌کنند.
\end{itemize}

\begin{center}\rule{0.5\linewidth}{0.5pt}\end{center}

\subsection{جزوه‌ی درس}

\paragraph{به زبان ساده}

اگر همان آزمایش را دوباره روی همان آدم اجرا کنید، داده‌ی کمی متفاوت و در
نتیجه برآورد پارامتر کمی متفاوتی می‌گیرید. \textbf{خطای استاندارد
\lr{(SE)}} می‌گوید برآورد بین چنین تکرارهایی چقدر بالا و پایین می‌پرد؛
\textbf{بازه‌ی اطمینان} محدوده‌ای است که از روی آن ساخته می‌شود (تقریباً
برآورد \ensuremath{\pm} 2 \lr{SE).} نمی‌توانیم آزمایش را بارها تکرار کنیم، پس این پرش را
به سه روش تخمین می‌زنیم: (1) از اینکه قله‌ی درست‌نمایی چقدر تیز است (قله‌ی
تیز یعنی داده پارامتر را دقیق تعیین می‌کند؛ این روش \textbf{\lr{Wald}} با
استفاده از \textbf{\lr{Hessian}} است)، (2) با دنبال کردن اینکه وقتی یک
پارامتر را جابه‌جا می‌کنیم برازش با چه سرعتی بدتر می‌شود
(\textbf{\lr{profile likelihood}})، و (3) با شبیه‌سازی تعداد زیادی آزمایش
ساختگی از مدل برازش‌شده و برازش دوباره‌ی هرکدام (\textbf{\lr{bootstrap}}).
نیمه‌ی دوم هفته درباره‌ی یک دام هشدار می‌دهد: میانگین گرفتن داده‌ها روی
افراد می‌تواند شکلی از منحنی بسازد که \textbf{هیچ فردی} آن را ندارد. اگر
هر یادگیرنده به‌طور نمایی ولی با سرعتی متفاوت بهتر شود، میانگین گروه شبیه
قانون توانی می‌شود.

\textbf{زمان‌بندی پیشنهادی (60 دقیقه):} توزیع‌های نمونه‌گیری 10 · خطای
استاندارد از \lr{Hessian} 15 · \lr{profile likelihood} 8 ·
\lr{bootstrap} 10 · چند شرکت‌کننده و میانگین‌گیری 17.

\subsubsection{1. برآوردها توزیع نمونه‌گیری دارند}

اگر همان آزمایش را روی همان شرکت‌کننده (با همان
\(\boldsymbol\theta^\ast\) واقعی) تکرار می‌کردیم، داده‌ی متفاوت و در نتیجه
\(\hat{\boldsymbol\theta}\) متفاوتی می‌گرفتیم. توزیع
\(\hat{\boldsymbol\theta}\) روی تکرارهای فرضی، \textbf{توزیع نمونه‌گیری}
آن است.

\begin{itemize}
\tightlist
\item
  \textbf{خطای استاندارد} \lr{(SE)} یک پارامتر، انحراف معیار توزیع
  نمونه‌گیری‌اش است: اینکه \emph{برآورد} بین تکرارها چقدر تغییر می‌کند.
\item
  یک \textbf{بازه‌ی اطمینان 95\%} با رویه‌ای ساخته می‌شود که در تکرارها، در
  95\% موارد \(\theta^\ast\) را در بر می‌گیرد. (این \emph{نه} به معنای
  «احتمال 95\% که \(\theta^\ast\) در همین بازه‌ی مشخص باشد» است؛ آن یک
  \lr{credible interval} بیزی است، هفته‌ی 9.)
\end{itemize}

نمی‌توانیم آزمایش را بارها تکرار کنیم، پس توزیع نمونه‌گیری را به سه روش
تقریب می‌زنیم.

\subsubsection{\texorpdfstring{2. خطای استاندارد از انحنای درست‌نمایی
\lr{(Wald)}}{2. خطای استاندارد از انحنای درست‌نمایی (Wald)}}

\textbf{شهود.} اگر \lr{NLL} یک کمینه‌ی تیز و باریک داشته باشد، تغییرات
کوچک در \(\boldsymbol\theta\) داده‌ها را خیلی کم‌احتمال‌تر می‌کنند، پس
داده‌ها \(\boldsymbol\theta\) را دقیق تعیین می‌کنند. یک کمینه‌ی پهن و تخت
یعنی مقادیر زیادی تقریباً به یک اندازه خوب برازش می‌شوند، پس عدم قطعیت
زیاد است. انحنا با مشتق‌های دوم اندازه‌گیری می‌شود.

\textbf{صورت‌بندی.} \textbf{\lr{Hessian}} تابع \lr{NLL} در نقطه‌ی \lr{MLE}
ماتریس مشتق‌های دوم است،

\[
\mathbf{H}_{mn} = \frac{\partial^2\, \mathrm{NLL}(\boldsymbol\theta)}{\partial\theta_m\,\partial\theta_n}\bigg|_{\boldsymbol\theta = \hat{\boldsymbol\theta}},
\]

که به آن \textbf{اطلاع فیشر مشاهده‌شده} \lr{(observed Fisher
information)} هم می‌گویند. برای نمونه‌های بزرگ:

\[
\widehat{\mathrm{Cov}}(\hat{\boldsymbol\theta}) \approx \mathbf{H}^{-1}, \qquad
\mathrm{SE}(\hat\theta_m) = \sqrt{[\mathbf{H}^{-1}]_{mm}}, \qquad
\text{\lr{95\% Wald CI: }} \hat\theta_m \pm 1.96\,\mathrm{SE}(\hat\theta_m).
\]

عناصر غیرقطری \(\mathbf{H}^{-1}\)، وقتی به همبستگی تبدیل شوند،
\textbf{\lr{trade-off} پارامترها} را نشان می‌دهند: همبستگی 0.9\ensuremath{-} بین دو
برآورد یعنی خطاهایشان تا حد زیادی همدیگر را جبران می‌کنند (همان دره‌ی مورب
هفته‌ی 3).

\textbf{بررسی درستی (دوجمله‌ای).} برای یک نسبت،
\(\mathrm{NLL}''(\theta) = y/\theta^2 + (n-y)/(1-\theta)^2\). در
\(\hat\theta = y/n\) این برابر \(n / [\hat\theta(1-\hat\theta)]\) است،
پس \(\mathrm{SE} = \sqrt{\hat\theta(1-\hat\theta)/n}\)؛ همان فرمول آشنای
کتاب‌های درسی.

\textbf{در عمل:} \(\mathbf{H}\) را عددی با تفاضل‌های متناهی حساب کنید
(بخش \lr{A lab).} \lr{\nolinkurl{hess_inv}} که \lr{BFGS} برمی‌گرداند فقط
یک تقریب خام است، پس برای \lr{SE} از آن استفاده نکنید.

\textbf{وقتی \lr{Wald} جواب نمی‌دهد:} نمونه‌های کوچک، برآوردهای نزدیک یک
کران (مثلاً سطح یادآوری \(\hat A\) نزدیک 1)، درست‌نمایی‌هایی که شدیداً
غیرمربعی‌اند، یا پارامترهایی با توزیع نمونه‌گیری چوله (نرخ‌ها، واریانس‌ها).
در این حالت‌ها بازه‌ی متقارن \(\pm 1.96\) \lr{SE} ممکن است تا مقادیر
ناممکن کشیده شود.

\begin{quote}
\textbf{پیشرفته: چرا این کار می‌کند.} نزدیک کمینه، هر \lr{NLL} همواری
تقریباً مربعی است (بسط تیلور مرتبه‌ی دوم). یک \lr{NLL} مربعی دقیقاً
\lr{NLL} یک توزیع نرمال چندمتغیره است که معکوس کوواریانسش همان
\lr{Hessian} است. همراه با نرمال بودن مجانبی \lr{MLE} (هفته‌ی 4)، این
نتیجه را می‌دهد. همین تقریب مربعی، به نام \emph{تقریب \lr{Laplace}}، در
مقایسه‌ی بیزی مدل‌ها (هفته‌ی 10) دوباره ظاهر می‌شود.
\end{quote}

\paragraph{خواندن ریاضی}

\begin{itemize}
\tightlist
\item
  \(\mathbf{H}_{mn} = \partial^2 \mathrm{NLL}/\partial\theta_m\partial\theta_n\)
  در نقطه‌ی \(\hat{\boldsymbol\theta}\): ماتریس مشتق‌های دوم، یعنی
  \textbf{انحنای} کاسه‌ی \lr{NLL} در کف آن. درایه‌های قطری: کاسه در امتداد
  هر پارامتر چقدر شیب‌دار است. درایه‌های غیرقطری: آیا کاسه کج است
  (پارامترها \lr{trade-off} دارند).
\item
  \(\widehat{\mathrm{Cov}}(\hat{\boldsymbol\theta}) \approx \mathbf{H}^{-1}\):
  کاسه‌ی خیلی خمیده (\(\mathbf{H}\) بزرگ) کوواریانس کوچک می‌دهد، یعنی
  برآورد دقیق. معکوس گرفتن همان چیزی است که «شیب‌داری» را به «عدم قطعیت»
  تبدیل می‌کند.
\item
  \(\mathrm{SE}(\hat\theta_m) = \sqrt{[\mathbf{H}^{-1}]_{mm}}\): درایه‌ی
  قطری \(m\)ام ماتریس معکوس را بردارید و از آن جذر بگیرید.
\item
  \textbf{بررسی با یک پارامتر:} برای یک نسبت،
  \(\mathrm{NLL}'' = n/[\hat\theta(1-\hat\theta)]\)، پس
  \(\mathrm{SE} = \sqrt{\hat\theta(1-\hat\theta)/n}\). با 30 از 50:
  \(\sqrt{0.6 \times 0.4/50} = 0.069\) و بازه‌ی 95\% برابر
  \(0.60 \pm 1.96 \times 0.069 = [0.46,\allowbreak 0.74]\) است.
\item
  \textbf{همبستگی از کوواریانس:}
  \(\mathrm{corr} = \mathrm{Cov}_{12}/(\mathrm{SE}_1\mathrm{SE}_2)\).
  مقداری مثل \ensuremath{-}0.9 یعنی دو برآورد به‌شدت \lr{trade-off} دارند.
\end{itemize}

\subsubsection{\texorpdfstring{3. بازه‌های \lr{profile
likelihood}}{3. بازه‌های profile likelihood}}

برای پارامتر مورد نظر \(\psi\) (مثلاً \(d\))، \textbf{\lr{NLL profile}}
با ثابت گرفتن \(\psi\) روی یک شبکه و کمینه کردن \lr{NLL} روی همه‌ی
پارامترهای \emph{دیگر} در هر نقطه‌ی شبکه به دست می‌آید:

\[
\mathrm{NLL}_{\text{\lr{prof}}}(\psi) = \min_{\text{\lr{other params}}} \mathrm{NLL}(\psi, \text{\lr{others}}).
\]

\textbf{بازه‌ی نسبت درست‌نمایی \lr{(profile)}} 95\%، مجموعه‌ی مقادیری از
\(\psi\) است که خیلی بدتر از بهترین نیستند:

\[
\big\{\psi : 2\,[\mathrm{NLL}_{\text{\lr{prof}}}(\psi) - \mathrm{NLL}_{\min}] \le 3.84\big\}, \qquad 3.84 = \chi^2_{1,\,0.95}.
\]

بازه‌های \lr{profile} شکل واقعی درست‌نمایی را دنبال می‌کنند، پس می‌توانند
\textbf{نامتقارن} باشند و کران‌های طبیعی را رعایت کنند. در نمونه‌های کوچک
معمولاً از بازه‌های \lr{Wald} دقیق‌ترند.

\paragraph{خواندن ریاضی}

\begin{itemize}
\tightlist
\item
  \(\mathrm{NLL}_{\text{\lr{prof}}}(\psi)\): پارامتر مورد نظر \(\psi\) را روی
  یک مقدار ثابت کنید، بگذارید بهینه‌ساز بهترین مقدار همه‌ی پارامترهای
  \emph{دیگر} را انتخاب کند، و \lr{NLL} حاصل را ثبت کنید. برای مقدارهای
  زیادی از \(\psi\) تکرار کنید. نتیجه یک منحنی یک‌بعدی است.
\item
  بازه هر \(\psi\)ای را نگه می‌دارد که \lr{NLL} نیم‌رخ آن حداکثر
  \(3.84/2 = 1.92\) واحد بالاتر از کمینه باشد. عدد 3.84 نقطه‌ی 95\% توزیع
  \(\chi^2\) با 1 درجه‌ی آزادی است و از آزمون نسبت درست‌نمایی می‌آید (هفته‌ی
  7).
\item
  برخلاف \lr{Wald}، این منحنی مجبور نیست سهمی متقارن باشد، پس بازه
  می‌تواند نامتقارن باشد و هیچ‌وقت تا مقدارهای ناممکن کشیده نمی‌شود.
\end{itemize}

\subsubsection{\texorpdfstring{4. \lr{bootstrap}
پارامتری}{4. bootstrap پارامتری}}

جوابی شبیه‌سازی‌محور به «اگر آزمایش را تکرار می‌کردیم چه؟»:

\begin{enumerate}
\def\labelenumi{\arabic{enumi}.}
\tightlist
\item
  مدل را روی داده برازش کنید: \(\hat{\boldsymbol\theta}\).
\item
  \(B\) مجموعه‌داده‌ی جدید از مدل با \(\hat{\boldsymbol\theta}\) و با همان
  طرح آزمایش واقعی شبیه‌سازی کنید.
\item
  هر مجموعه‌داده را دوباره برازش کنید:
  \(\hat{\boldsymbol\theta}^{*(1)},\allowbreak \dots,\allowbreak \hat{\boldsymbol\theta}^{*(B)}\).
\item
  انحراف معیار برآوردهای \lr{bootstrap} تقریب \lr{SE} است. صدک‌های 2.5 و
  97.5 یک \textbf{بازه‌ی اطمینان صدکی} \lr{(percentile CI)} می‌دهند.
\end{enumerate}

\lr{bootstrap} به مشتق نیاز ندارد، برای هر آماره‌ای کار می‌کند (مثلاً
نیمه‌عمر یا تفاوت دو شرط)، و شکل کامل و همبستگی‌های توزیع نمونه‌گیری را
نشان می‌دهد. هزینه‌اش \(B\) بار برازش است (\(B \approx 500\) تا \(2000\)).
\lr{bootstrap} \textbf{ناپارامتری} به‌جای شبیه‌سازی از مدل، \lr{trial}های
مشاهده‌شده را با جایگذاری بازنمونه‌گیری می‌کند.

\textbf{مقایسه:}

{\def\LTcaptype{none} % do not increment counter
\begin{longtable}[]{@{}
  >{\raggedright\arraybackslash}p{(\linewidth - 8\tabcolsep) * \real{0.2000}}
  >{\raggedright\arraybackslash}p{(\linewidth - 8\tabcolsep) * \real{0.2000}}
  >{\raggedright\arraybackslash}p{(\linewidth - 8\tabcolsep) * \real{0.2000}}
  >{\raggedright\arraybackslash}p{(\linewidth - 8\tabcolsep) * \real{0.2000}}
  >{\raggedright\arraybackslash}p{(\linewidth - 8\tabcolsep) * \real{0.2000}}@{}}
\toprule\noalign{}
\begin{minipage}[b]{\linewidth}\raggedright
روش
\end{minipage} & \begin{minipage}[b]{\linewidth}\raggedright
هزینه
\end{minipage} & \begin{minipage}[b]{\linewidth}\raggedright
متقارن؟
\end{minipage} & \begin{minipage}[b]{\linewidth}\raggedright
کران را رعایت می‌کند؟
\end{minipage} & \begin{minipage}[b]{\linewidth}\raggedright
نیاز دارد به
\end{minipage} \\
\midrule\noalign{}
\endhead
\bottomrule\noalign{}
\endlastfoot
\lr{Wald (Hessian)} & یک \lr{Hessian} & همیشه & نه & \lr{NLL} هموار،
\(N\) بزرگ \\
\lr{profile likelihood} & یک برازش برای هر نقطه‌ی شبکه & می‌تواند نامتقارن
باشد & بله & یک درست‌نمایی \\
\lr{bootstrap} پارامتری & \(B\) برازش & می‌تواند نامتقارن باشد & بله &
مدلی که بتوانید شبیه‌سازی کنید \\
\end{longtable}
}

\subsubsection{5. برازش برای چند شرکت‌کننده}

آزمایش‌های واقعی \(J\) شرکت‌کننده دارند (اندیس \(j\)). سه راهبرد کلی وجود
دارد:

{\def\LTcaptype{none} % do not increment counter
\begin{longtable}[]{@{}
  >{\raggedright\arraybackslash}p{(\linewidth - 6\tabcolsep) * \real{0.2500}}
  >{\raggedright\arraybackslash}p{(\linewidth - 6\tabcolsep) * \real{0.2500}}
  >{\raggedright\arraybackslash}p{(\linewidth - 6\tabcolsep) * \real{0.2500}}
  >{\raggedright\arraybackslash}p{(\linewidth - 6\tabcolsep) * \real{0.2500}}@{}}
\toprule\noalign{}
\begin{minipage}[b]{\linewidth}\raggedright
راهبرد
\end{minipage} & \begin{minipage}[b]{\linewidth}\raggedright
چه چیزی برازش می‌شود
\end{minipage} & \begin{minipage}[b]{\linewidth}\raggedright
مزایا
\end{minipage} & \begin{minipage}[b]{\linewidth}\raggedright
معایب
\end{minipage} \\
\midrule\noalign{}
\endhead
\bottomrule\noalign{}
\endlastfoot
\textbf{تجمیعی} \lr{(complete pooling)} & یک مدل روی داده‌های
میانگین‌گیری‌شده & پایدار؛ با \lr{trial}های کم برای هر نفر هم کار می‌کند &
فرض می‌کند همه یکسان‌اند؛ \textbf{می‌تواند شکل تابعی را تحریف کند} \\
\textbf{فردی} \lr{(no pooling)} & یک مدل برای هر شرکت‌کننده؛ بعد تحلیل
\(J\) مجموعه برآورد & تفاوت‌های فردی را رعایت می‌کند؛ روش استاندارد & با
\lr{trial} کم برآوردها پرنویزند؛ تحلیل مرحله‌ی دوم عدم قطعیت برآورد را
نادیده می‌گیرد \\
\textbf{سلسله‌مراتبی} \lr{(partial pooling)} & پارامترهای فردی از یک
توزیع گروهی، که با هم برآورد می‌شوند & بهترینِ هر دو؛ \lr{shrinkage} اصولی
& پیچیده‌تر؛ معمولاً بیزی (هفته‌ی 13) \\
\end{longtable}
}

\textbf{خطاهای ناشی از میانگین‌گیری.} میانگین منحنی‌های غیرخطی معمولاً
منحنی‌ای از همان خانواده \textbf{نیست}. مثلاً میانگین منحنی‌های نمایی با
نرخ‌های متفاوت نمایی نمی‌شود: یادگیرنده‌های سریع زود به پایان می‌رسند، پس
نرخ نسبی تغییر میانگین با گذشت زمان \emph{کند می‌شود}؛ و این دقیقاً ویژگی
تعریف‌کننده‌ی تابع توانی است (هفته‌ی 2، نرخ نسبی فراموشی).

\begin{quote}
\textbf{پیشرفته: یک نتیجه‌ی دقیق.} فرض کنید هر شرکت‌کننده از
\(e^{-\lambda_j t}\) پیروی کند و نرخ‌ها بین افراد توزیع \lr{Gamma} داشته
باشند،
\(\lambda_j \sim \mathrm{Gamma}(\text{\lr{shape }} s, \text{\lr{rate }} b)\).
آن‌وقت منحنی میانگین گروه برابر است با
\[\mathbb{E}\big[e^{-\lambda t}\big] = \Big(1 + \frac{t}{b}\Big)^{-s},\]
که \emph{دقیقاً} یک تابع توانی جابه‌جاشده است (این تابع مولد گشتاور توزیع
\lr{Gamma} در نقطه‌ی \(-t\) است). همه به‌طور نمایی فراموش می‌کنند، ولی
میانگین با یک قانون توانی با نمای برابر با پارامتر شکل \lr{Gamma} یعنی
\(s\) فراموش می‌کند. این را در تمرین 4 عددی بررسی می‌کنید.
\end{quote}

\textbf{راهنمای عملی.} هر وقت هر شرکت‌کننده داده‌ی کافی برای مقید کردن مدل
دارد، افراد را برازش کنید؛ داده‌ی تجمیعی را فقط به‌عنوان خلاصه به کار
ببرید، هرگز به‌عنوان شاهدی درباره‌ی \emph{شکل} یک فرایند در سطح فرد. وقتی
داده‌های فردی کم است، از مدل‌های سلسله‌مراتبی استفاده کنید (هفته‌ی 13).

\paragraph{برای یک مهندس}

\begin{itemize}
\tightlist
\item
  \textbf{\lr{SE} روش \lr{Wald} از \lr{Hessian}} = معکوس اطلاعات فیشر،
  یعنی همان \textbf{\lr{CRLB}} که در نقطه‌ی برآورد محاسبه شده. \lr{MLE}
  به‌طور مجانبی به این کران می‌رسد، پس برای \(N\) بزرگ بازه‌ی \lr{Wald}
  همان بازه‌ی مبتنی بر \lr{CRLB} است.
\item
  \textbf{\lr{bootstrap} پارامتری} همان ارزیابی کارایی با \lr{Monte
  Carlo} است که برای یک تخمین‌گر یا گیرنده انجام می‌دهید (مثلاً شبیه‌سازی
  تعداد زیادی تحقق نویز برای به دست آوردن واریانس تخمین‌گر یا منحنی
  \lr{BER)}، فقط در مقدارهای برازش‌شده‌ی پارامتر.
\item
  \textbf{\lr{profile likelihood}} همان کاری است که با \emph{پارامترهای
  مزاحم \lr{(nuisance parameters)}} می‌کنید: آن‌ها را بهینه کنید و کنار
  بگذارید و به درست‌نمایی پارامتری که برایتان مهم است نگاه کنید.
\item
  \textbf{خطای ناشی از میانگین‌گیری} = تفاوت \textbf{میانگین آنسامبل} و
  \textbf{یک تحقق تنها}. میانگین گرفتن از افت‌های نمایی با ثابت‌های زمانی
  مختلف یک افت کندتر با دنباله‌ی سنگین می‌دهد؛ همان‌طور که میانگین تعداد
  زیادی تحقق کانال شبیه هیچ تحقق تکی نیست. فرایند «بین افراد ارگودیک»
  نیست: میانگین گروه، فرد را توصیف نمی‌کند.
\item
  \textbf{\lr{partial pooling}} (پیش‌نمایش هفته‌ی 13) شبیه تخمین \lr{SNR}
  هر کاربر با ترکیب چند اندازه‌گیری خودش و توزیع \lr{SNR} در کل سلول است:
  یک برآورد \lr{empirical Bayes.}
\end{itemize}

\paragraph{برای مدرس}

\begin{itemize}
\tightlist
\item
  \textbf{با سؤال «آیا 0.42 با 0.5 فرق دارد؟» شروع کنید} (بخش مقدمه).
  دانشجوها جواب بله/نه می‌خواهند؛ نشان دهید جواب کاملاً به \lr{SE} بستگی
  دارد.
\item
  \textbf{سه روش را کنار هم نشان دهید} روی همان داده (جدول بخش \lr{C
  lab).} پیام این نیست که «کدام درست است»، بلکه «وقتی درست‌نمایی خوش‌رفتار
  است با هم توافق دارند و نزدیک کران‌ها یا با داده‌ی کم فرق می‌کنند».
\item
  \textbf{نمایش میانگین‌گیری نقطه‌ی اوج جلسه است.} 5 منحنی نمایی فردی و
  میانگین آن‌ها را روی نمودار لگاریتمی\lr{--}لگاریتمی رسم کنید: میانگین تقریباً
  یک خط راست است (قانون توانی) در حالی که هر منحنی فردی خمیده است. از
  کلاس بپرسید پژوهشگری که فقط میانگین را دیده چه نتیجه‌ای می‌گیرد.
\item
  \textbf{سؤال‌های رایج.}

  \begin{itemize}
  \tightlist
  \item
    \emph{«فرق \lr{SD} و \lr{SE} چیست؟»} \lr{SD} پراکندگی \emph{داده‌ها}
    را توصیف می‌کند؛ \lr{SE} پراکندگی یک \emph{برآورد} در تکرارهای آزمایش
    را. \lr{SE} با داده‌ی بیشتر کم می‌شود؛ \lr{SD} نه.
  \item
    \emph{«آیا یک بازه‌ی اطمینان 95\% با احتمال 0.95 مقدار واقعی را
    دارد؟»} این بازه‌ی مشخص نه؛ 95\% بازه‌هایی که به این روش ساخته می‌شوند
    آن را دارند. \lr{(credible interval} بیزی هفته‌ی 9 همین تفسیر مستقیم
    را دارد.)
  \item
    \emph{«آیا همیشه باید افراد را جداگانه برازش کنم؟»} بله، اگر هر نفر
    داده‌ی کافی دارد؛ وگرنه از مدل‌های سلسله‌مراتبی استفاده کنید (هفته‌ی
    13).
  \end{itemize}
\end{itemize}

\begin{center}\rule{0.5\linewidth}{0.5pt}\end{center}

\subsection{مثال حل‌شده}

\subsubsection{بخش 1: سه بازه برای یک احتمال یادآوری}

یک شرکت‌کننده \(y = 30\) از \(n = 50\) آیتم را به یاد می‌آورد؛
\(\hat\theta = 0.60\).

\textbf{\lr{Wald.}}
\(\mathrm{SE} = \sqrt{0.6 \times 0.4 / 50} = 0.0693\)، پس بازه‌ی اطمینان
95\% برابر \(0.60 \pm 1.96 \times 0.0693 = [0.464,\allowbreak 0.736]\) است.

\textbf{\lr{profile} (نسبت درست‌نمایی).} معادله‌ی
\(2[\ell(0.6) - \ell(\theta)] = 3.84\) را در دو طرف 0.6 عددی حل کنید.
بازه‌ی 95\% برابر \([0.462, 0.728]\) است که کمی \textbf{نامتقارن} است:
0.138 پایین‌تر از برآورد و 0.128 بالاتر از آن می‌رود، چون درست‌نمایی به سمت
0.5، جایی که واریانس دوجمله‌ای بیشینه است، چوله است.

\textbf{تفسیر.} با 50 آیتم، یک احتمال یادآوری فقط تا حدود \ensuremath{\pm}0.13 معلوم
است. تفاوت 0.10 بین دو شرط با این اندازه به‌تنهایی قانع‌کننده نیست.

\subsubsection{بخش 2: میانگین‌گیری از دو یادگیرنده‌ی نمایی}

شرکت‌کننده‌ی 1 با \(\lambda = 0.1\) و شرکت‌کننده‌ی 2 با \(\lambda = 1.0\)
فراموش می‌کند. منحنی میانگین \(\bar f(t) = (e^{-0.1t} + e^{-t})/2\):

{\def\LTcaptype{none} % do not increment counter
\begin{longtable}[]{@{}
  >{\raggedright\arraybackslash}p{(\linewidth - 10\tabcolsep) * \real{0.1667}}
  >{\raggedright\arraybackslash}p{(\linewidth - 10\tabcolsep) * \real{0.1667}}
  >{\raggedright\arraybackslash}p{(\linewidth - 10\tabcolsep) * \real{0.1667}}
  >{\raggedright\arraybackslash}p{(\linewidth - 10\tabcolsep) * \real{0.1667}}
  >{\raggedright\arraybackslash}p{(\linewidth - 10\tabcolsep) * \real{0.1667}}
  >{\raggedright\arraybackslash}p{(\linewidth - 10\tabcolsep) * \real{0.1667}}@{}}
\toprule\noalign{}
\begin{minipage}[b]{\linewidth}\raggedright
\(t\)
\end{minipage} & \begin{minipage}[b]{\linewidth}\raggedright
0
\end{minipage} & \begin{minipage}[b]{\linewidth}\raggedright
1
\end{minipage} & \begin{minipage}[b]{\linewidth}\raggedright
2
\end{minipage} & \begin{minipage}[b]{\linewidth}\raggedright
5
\end{minipage} & \begin{minipage}[b]{\linewidth}\raggedright
10
\end{minipage} \\
\midrule\noalign{}
\endhead
\bottomrule\noalign{}
\endlastfoot
\(\bar f(t)\) & 1.000 & 0.636 & 0.477 & 0.307 & 0.184 \\
نرخ نسبی فراموشی \(-\bar f'(t)/\bar f(t)\) & 0.55 & 0.36 & 0.23 & 0.11 &
0.10 \\
\end{longtable}
}

هر شرکت‌کننده یک نرخ نسبی \emph{ثابت} دارد (0.1 و 1.0)، ولی نرخ میانگین
از 0.55 به 0.10 \textbf{کاهش} می‌یابد؛ نشانه‌ی یک منحنی شبه‌توانی. دلیل: در
زمان‌های بعدی، فراموش‌کننده‌ی کند که نرخش پایین است بر میانگین غالب است.
\textbf{تفسیر:} پژوهشگری که فقط میانگین را برازش کند نتیجه می‌گیرد که
«خاطره‌ها با بالا رفتن سن مقاوم‌تر می‌شوند»؛ ادعایی که درباره‌ی هیچ‌کدام از
دو شرکت‌کننده درست نیست.

\begin{center}\rule{0.5\linewidth}{0.5pt}\end{center}

\subsection{\texorpdfstring{\lr{lab} (2 ساعت)}{lab (2 ساعت)}}

\subsubsection{اهداف}

\begin{enumerate}
\def\labelenumi{\arabic{enumi}.}
\tightlist
\item
  محاسبه‌ی بازه‌های \lr{Wald}، \lr{profile} و \lr{bootstrap} برای مدل
  توانی برازش‌شده روی داده‌های هفته‌ی 2 و مقایسه‌ی آن‌ها.
\item
  نمایش \lr{trade-off} پارامترها از طریق همبستگی برآوردشده و ابر نقاط
  \lr{bootstrap.}
\item
  شبیه‌سازی 20 یادگیرنده‌ی نمایی، برازش مدل‌های نمایی و توانی روی
  \textbf{میانگین} و روی \textbf{هر فرد}، و تفسیر تناقض.
\end{enumerate}

\textbf{زمان‌بندی:} بخش \lr{A} 20 دقیقه · بخش \lr{B} 15 دقیقه · بخش
\lr{C} 25 دقیقه · بخش \lr{D} 40 دقیقه · بخش \lr{E} 20 دقیقه.

\subsubsection{آماده‌سازی}

\begin{quote}
\textbf{\lr{Code} · کد:} \lr{\nolinkurl{code/week05.py}}, \lr{cell} 1 ·
سلول 1
\end{quote}

\subsubsection{\texorpdfstring{بخش \lr{A:} خطای استاندارد \lr{Wald} از
یک \lr{Hessian} عددی (20
دقیقه)}{بخش A: خطای استاندارد Wald از یک Hessian عددی (20 دقیقه)}}

\begin{quote}
\textbf{\lr{Code} · کد:} \lr{\nolinkurl{code/week05.py}}, \lr{cell} 2 ·
سلول 2
\end{quote}

\textbf{خروجی مورد انتظار:}

\Needspace{9\baselineskip}
\begin{latin}
\begin{Verbatim}
   estimate      SE   lower   upper
A    0.8773  0.0626  0.7546  1.0000
d    0.4214  0.0532  0.3171  0.5257
Estimated correlation between A-hat and d-hat: 0.751
\end{Verbatim}
\end{latin}

\textbf{تفسیر:} هر دو بازه‌ی 95\% مقادیر واقعی (\(A^\ast = 0.9\)،
\(d^\ast = 0.5\)) را در بر دارند، پس «خطایی» که در هفته‌های 3\lr{--}4 نگرانش
بودیم کاملاً در حد خطای نمونه‌گیری است. دقت کنید که بازه‌ی \lr{Wald} برای
\(A\) اتفاقاً درست روی 1.000، یعنی مرز مقادیر ممکن، تمام می‌شود. با داده‌ای
کمی پرنویزتر، بازه‌ی متقارن \lr{Wald} مقادیر ناممکن (احتمال بیشتر از 1)
را هم شامل می‌شد، در حالی که بازه‌ی \lr{bootstrap} نمی‌تواند. همبستگی مثبت
همان دره‌ی \lr{trade-off} هفته‌ی 3 است: وقتی برآورد \(A\) زیادی بالاست،
برآورد \(d\) هم معمولاً زیادی بالاست.

\subsubsection{\texorpdfstring{بخش \lr{B:} بازه‌ی \lr{profile likelihood}
برای \(d\) (15
دقیقه)}{بخش B: بازه‌ی profile likelihood برای d (15 دقیقه)}}

\begin{quote}
\textbf{\lr{Code} · کد:} \lr{\nolinkurl{code/week05.py}}, \lr{cell} 3 ·
سلول 3
\end{quote}

\textbf{خروجی مورد انتظار:}

\Needspace{6\baselineskip}
\begin{latin}
\begin{Verbatim}
Profile 95% CI for d: [0.320, 0.525]  (grid step 0.005)
\end{Verbatim}
\end{latin}

\subsubsection{\texorpdfstring{بخش \lr{C: bootstrap} پارامتری (25
دقیقه)}{بخش C: bootstrap پارامتری (25 دقیقه)}}

\begin{quote}
\textbf{\lr{Code} · کد:} \lr{\nolinkurl{code/week05.py}}, \lr{cell} 4 ·
سلول 4
\end{quote}

\textbf{خروجی مورد انتظار} (حدود 30 تا 60 ثانیه طول می‌کشد):

\Needspace{9\baselineskip}
\begin{latin}
\begin{Verbatim}
   boot SE    2.5%  97.5%
A   0.0563  0.7724  1.000
d   0.0501  0.3304  0.528
Bootstrap correlation: 0.703
\end{Verbatim}
\end{latin}

\textbf{مقایسه کنید:} بازه‌های \lr{Wald}، \lr{profile} و \lr{bootstrap}
برای \(d\) را کنار هم بگذارید. شبیه هم‌اند؟ کدام نامتقارن است؟ با این
داده‌ها (8 خانه × 50 آیتم) هر سه روش باید نسبتاً با هم توافق داشته باشند.
با آیتم‌های کمتر یا برآوردهای نزدیک کران بیشتر از هم فاصله می‌گیرند (تمرین
5).

\subsubsection{\texorpdfstring{بخش \lr{D:} قانون توانی تمرین به‌عنوان
خطای میانگین‌گیری (40
دقیقه)}{بخش D: قانون توانی تمرین به‌عنوان خطای میانگین‌گیری (40 دقیقه)}}

\textbf{طرح.} \(J = 20\) شرکت‌کننده هرکدام 60 بلوک تمرین در یک تکلیف
جستجوی دیداری انجام می‌دهند. هر شرکت‌کننده به‌طور \textbf{نمایی} بهتر
می‌شود:

\[
\mathrm{RT}_j(m) = E + B\,e^{-r_j (m-1)} + \varepsilon, \qquad \varepsilon \sim \mathcal{N}(0, 40 \text{\lr{ ms}}),
\]

با مجانب \(E = 400\) \lr{ms}، بهبود اولیه‌ی \(B = 600\) \lr{ms} و نرخ‌های
یادگیری فردی
\(r_j \sim \mathrm{Gamma}(\text{\lr{shape}} = 0.7, \text{\lr{scale}} = 0.5)\)
(میانگین 0.35 در هر بلوک؛ توزیعی چوله با چند یادگیرنده‌ی خیلی سریع و چند
یادگیرنده‌ی خیلی کند). اینجا \(m\) شماره‌ی بلوک تمرین و \(\varepsilon\)
نویز میانگین زمان واکنش هر بلوک است.

\begin{quote}
\textbf{\lr{Code} · کد:} \lr{\nolinkurl{code/week05.py}}, \lr{cell} 5 ·
سلول 5
\end{quote}

\textbf{خروجی مورد انتظار:}

\Needspace{10\baselineskip}
\begin{latin}
\begin{Verbatim}
True learning rates: [1.000e-03 7.000e-03 7.000e-03 2.000e-02 3.000e-02 4.500e-02 1.220e-01
 1.490e-01 1.720e-01 2.360e-01 2.820e-01 3.260e-01 3.870e-01 3.950e-01
 5.340e-01 7.100e-01 7.900e-01 1.002e+00 1.869e+00 2.534e+00]
Average curve, exp: E = 507.5, B = 423.8, r = 0.205, RMSD = 19.26 ms
Average curve, pow: E = 419.2, B = 579.4, r = 0.529, RMSD = 10.16 ms
\end{Verbatim}
\end{latin}

\begin{quote}
\textbf{\lr{Code} · کد:} \lr{\nolinkurl{code/week05.py}}, \lr{cell} 6 ·
سلول 6
\end{quote}

\textbf{خروجی مورد انتظار (خط خلاصه):}

\Needspace{7\baselineskip}
\begin{latin}
\begin{Verbatim}
Which model fits each participant better? {'exp': 20}
(with SciPy 1.18: {'exp': 19, 'pow': 1}; one participant with a near-zero learning rate is fitted differently)
\end{Verbatim}
\end{latin}

\begin{quote}
\textbf{\lr{Code} · کد:} \lr{\nolinkurl{code/week05.py}}, \lr{cell} 7 ·
سلول 7
\end{quote}

\subsubsection{\texorpdfstring{بخش \lr{E:} تفسیر (20
دقیقه)}{بخش E: تفسیر (20 دقیقه)}}

\begin{quote}
\textbf{\lr{Code} · کد:} \lr{\nolinkurl{code/week05.py}}, \lr{cell} 8 ·
سلول 8
\end{quote}

\textbf{خروجی مورد انتظار:}

\Needspace{10\baselineskip}
\begin{latin}
\begin{Verbatim}
Mean of true individual rates:      0.481
Mean of estimated individual rates: 0.449 (SD 0.626)
Exponential rate fitted to the average curve: 0.205
Correlation of true and estimated individual rates: 0.971
(with SciPy 1.18: estimated mean 0.457 (SD 0.621), correlation 0.968, because of the same participant)
\end{Verbatim}
\end{latin}

\textbf{برای هرکدام 2 تا 3 جمله بنویسید:} 1. منحنی میانگین با تابع توانی
بهتر توصیف می‌شود \lr{(RMSD} حدود 10 در برابر 19 \lr{ms)}، ولی همه‌ی 20
فرد (یا بسته به نسخه‌ی \lr{SciPy}، 19 نفر) با تابع نمایی بهتر توصیف
می‌شوند. کدام گزاره درباره‌ی \emph{یادگیری} موجه است و کدام نیست؟ 2. چرا
نرخ نمایی برازش‌شده روی منحنی میانگین از میانگین نرخ‌های فردی کوچک‌تر است؟
(راهنمایی: در اواخر تمرین کدام شرکت‌کننده‌ها بر میانگین غالب‌اند؟) 3. اینجا
با 60 بلوک برای هر نفر، برازش‌های فردی نرخ‌های واقعی را خوب بازیابی
می‌کنند. با فقط 10 بلوک برای هر نفر چه تغییر می‌کرد و آن‌وقت چه می‌کردید؟
(پیش‌نمایش هفته‌ی 13.)

\subsubsection{تمرین‌ها (به ترتیب سختی)}

\textbf{تمرین 1 \ensuremath{\star}.} بازه‌های اطمینان 95\% \lr{Wald} را برای \lr{MLE} مدل
\emph{نمایی} روی داده‌های هفته‌ی 2 حساب کنید. آیا مقادیر معقولی در بر
دارند؟ بعد آن‌ها را در «برآورد» ساختگی
\lr{\nolinkurl{(A, lam) = (1.0, 0.107)}} از هشدار \lr{Nelder\lr{--}Mead}
کران‌دار هفته‌ی قبل حساب کنید. در نقطه‌ای که کمینه نیست چه بر سر
\lr{Hessian} می‌آید؟

\textbf{تمرین 2 \ensuremath{\star}.} \lr{bootstrap} را با \(B = 50, 200, 1000\) و
\lr{seed}های مختلف دوباره اجرا کنید. حدود صدکی چقدر پایدارند؟ برای یک
بازه‌ی 95\%، \(B\) چقدر باید بزرگ باشد؟

\textbf{تمرین 3 \ensuremath{\star}\ensuremath{\star} (یک کمیت مشتق‌شده).} با نمونه‌های \lr{bootstrap} بخش
\lr{C} یک بازه‌ی 95\% برای \textbf{نیمه‌عمر} مدل توانی،
\(t_{1/2} = 2^{1/d} - 1\)، حساب کنید. چرا این کار با \lr{bootstrap} آسان
و با \lr{Hessian} دست‌وپاگیر است؟

\textbf{تمرین 4 \ensuremath{\star}\ensuremath{\star} (نتیجه‌ی \lr{Gamma\lr{--}}توانی).} 100 000 نرخ
\(\lambda \sim \mathrm{Gamma}(\text{\lr{shape }} s = 0.8, \text{\lr{rate }} b = 2)\)
بکشید، \(e^{-\lambda t}\) را روی آن‌ها برای \(t \in [0, 30]\) میانگین
بگیرید و منحنی تحلیلی \((1 + t/b)^{-s}\) را روی آن رسم کنید (کادر
«پیشرفته» بخش 5 جزوه). بعد توابع نمایی و توانی را روی منحنی میانگین
برازش کنید.

\textbf{تمرین 5 \ensuremath{\star}\ensuremath{\star}\ensuremath{\star} (پیشرفته: پوشش).} بازه‌های 95\% واقعاً چند وقت یک بار
مقدار واقعی را در بر می‌گیرند؟ 300 مجموعه‌داده از مدل توانی (\(A = 0.9\)،
\(d = 0.5\)) با \(n = 10\) آیتم در هر فاصله شبیه‌سازی کنید و برای هرکدام
بازه‌های \lr{Wald} و \lr{profile} \(d\) را حساب کنید. پوشش (نسبت بازه‌هایی
که 0.5 را در بر دارند) هر روش را گزارش کنید. کدام به 95\% نزدیک‌تر است؟
با \(n = 50\) تکرار کنید.

\subsubsection{\texorpdfstring{\lr{lab check-off}}{lab check-off}}

نشان دهید: (1) جدولی که بازه‌های 95\% \lr{Wald}، \lr{profile} و
\lr{bootstrap} برای \(d\) را مقایسه می‌کند؛ (2) نمودار پراکنش
\lr{bootstrap} با یک جمله درباره‌ی \lr{trade-off}؛ (3) مقایسه‌ی فردی در
برابر میانگین و پاسخ شما به سؤال تفسیری 1.

\begin{center}\rule{0.5\linewidth}{0.5pt}\end{center}

\subsection{\texorpdfstring{ارتباط با هوش مصنوعی و
\lr{LLM}ها}{ارتباط با هوش مصنوعی و LLMها}}

\textbf{نمره‌های \lr{benchmark} هم خطای استاندارد دارند.} مدل‌های هوش
مصنوعی روی \lr{benchmark}ها مقایسه می‌شوند، و جدول‌های رتبه‌بندی
\lr{(leaderboard)} اغلب مدل‌هایی را رتبه‌بندی می‌کنند که یکی دو درصد با هم
فرق دارند. یک \lr{benchmark} یک \emph{نمونه} از آیتم‌هاست، پس دقتش، درست
مثل برآورد یک پارامتر، خطای نمونه‌گیری دارد. دو درس این هفته مستقیماً به
کار می‌آیند: 1. \textbf{عدم قطعیت را گزارش کنید.} \lr{bootstrap} روی
آیتم‌ها یک بازه‌ی اطمینان برای دقت و برای تفاوت دو مدل می‌دهد. 2.
\textbf{به طرح احترام بگذارید.} هر دو مدل به \emph{همان} آیتم‌ها جواب
داده‌اند، پس مقایسه جفتی است، مثل یک طرح درون‌آزمودنی. نمونه‌گیری مجدد
مستقل از دو مدل، دشواری مشترک آیتم‌ها را نادیده می‌گیرد و بازه‌هایی بیش از
حد پهن می‌دهد.

دشواری آیتم‌ها هم خیلی متفاوت است، که نسخه‌ی سطح آیتمِ ناهمگنی شرکت‌کننده‌ها
در بخش \lr{D} است: یک دقت میانگین واحد آن را پنهان می‌کند.

\begin{quote}
\textbf{\lr{Code} · کد:} \lr{\nolinkurl{code/week05.py}}, \lr{cell} 9 ·
سلول 9
\end{quote}

\textbf{خروجی مورد انتظار:}

\Needspace{8\baselineskip}
\begin{latin}
\begin{Verbatim}
Accuracy A = 0.74, B = 0.67, difference = 0.07
95% CI, paired bootstrap:   [-0.03, 0.17]
95% CI, unpaired bootstrap: [-0.06, 0.20]
\end{Verbatim}
\end{latin}

7 درصد برتری روی 100 آیتم به‌وضوح با صفر فرق ندارد. جدول رتبه‌بندی‌ای که
\lr{A} را بدون بازه بالای \lr{B} می‌گذارد، بیش از حد مطمئن است.

\textbf{امتحانش کنید (بدون نیاز به \lr{API).}} - \ensuremath{\star} اگر تفاوت واقعی ثابت
بماند، \lr{benchmark} به چند آیتم نیاز دارد تا بازه‌ی جفتی 0 را شامل
نشود؟ با شبیه‌سازی پیدا کنید. - \ensuremath{\star}\ensuremath{\star} \lr{SE} والد \lr{(Wald)} هر دقت،
\(\sqrt{p(1-p)/n}\)، را حساب کنید و با \lr{SE} حاصل از \lr{bootstrap}
مقایسه کنید. - \ensuremath{\star}\ensuremath{\star} بحث: مقاله‌ای «مدل \lr{A:} 74\%، مدل \lr{B:} 67\%» را
بدون بازه گزارش می‌کند و نتیجه می‌گیرد که \lr{A} «بهتر استدلال می‌کند».
فهرست کنید از نویسندگان چه می‌پرسید.

\begin{center}\rule{0.5\linewidth}{0.5pt}\end{center}

\subsection{اشتباهات رایج و برداشت‌های نادرست}

\begin{enumerate}
\def\labelenumi{\arabic{enumi}.}
\tightlist
\item
  \textbf{قاطی کردن \lr{SE} با \lr{SD.}} \lr{SD} پراکندگی \emph{داده‌ها}
  را توصیف می‌کند؛ \lr{SE} عدم قطعیت یک \emph{برآورد} را توصیف می‌کند و با
  بزرگ شدن داده‌ها کوچک می‌شود.
\item
  \textbf{«با احتمال 95\% مقدار واقعی در بازه‌ی اطمینان من است.»} در
  بازه‌ی اطمینان فراوانی‌گرا، 95\% به رفتار بلندمدت رویه برمی‌گردد. این
  گزاره‌ی احتمالاتی مال \lr{credible interval}های بیزی است (هفته‌ی 9).
\item
  \textbf{استفاده از \lr{\nolinkurl{hess_inv}} در \lr{BFGS} برای خطای
  استاندارد.} یک محصول جانبی خام بهینه‌ساز است. \lr{Hessian} را درست حساب
  کنید.
\item
  \textbf{اعتماد به بازه‌های \lr{Wald} نزدیک کران‌ها.} اگر \(\hat\theta\)
  روی کران یا نزدیک آن است، یا بازه از مقادیر ناممکن عبور می‌کند، از
  بازه‌های \lr{profile} یا \lr{bootstrap} استفاده کنید.
\item
  \textbf{گزارش پارامترها بدون عدم قطعیت.} مقایسه‌ی دو برآورد نقطه‌ای
  («شرکت‌کننده‌ی 1 دارای \(d = 0.42\) و شرکت‌کننده‌ی 2 دارای \(d = 0.51\)
  است») بدون بازه بی‌معنی است.
\item
  \textbf{استنباط فرایندهای سطح فرد از داده‌های میانگین‌گیری‌شده.} میانگین
  منحنی‌های غیرخطی شکل تابعی را عوض می‌کند (بخش \lr{D).} این برای آیتم‌ها،
  \lr{trial}ها و شرکت‌کننده‌ها یکسان صادق است.
\item
  \textbf{با برآوردهای فردی مثل داده رفتار کردن.} آزمون \lr{t} مرحله‌ی
  دوم روی پارامترهای برازش‌شده نادیده می‌گیرد که بعضی برآوردها خیلی
  پرنویزتر از بقیه‌اند. مدل‌های سلسله‌مراتبی این را درست می‌کنند (هفته‌ی 13).
\end{enumerate}

\begin{center}\rule{0.5\linewidth}{0.5pt}\end{center}

\subsection{خودآزمایی}

\textbf{1.} توزیع نمونه‌گیری یک برآوردگر چیست و چرا نمی‌توانیم مستقیماً آن
را ببینیم؟ \textbf{پاسخ} توزیع برآوردهایی که در تکرارهای فرضی بسیار زیاد
آزمایش با همان پارامترهای واقعی می‌گرفتیم. فقط یک مجموعه‌داده داریم، پس آن
را با \lr{Hessian} (نظریه‌ی مجانبی)، \lr{profile likelihood} یا شبیه‌سازی
\lr{(bootstrap)} تقریب می‌زنیم.

\textbf{2.} \lr{NLL} در یک جهت از فضای پارامتر کمینه‌ی خیلی تختی دارد.
این درباره‌ی خطاهای استاندارد و همبستگی برآوردها چه می‌گوید؟ \textbf{پاسخ}
انحنای کم یعنی درایه‌ی کوچک \lr{Hessian} در آن جهت، پس واریانس بزرگ در
\(\mathbf{H}^{-1}\): \lr{SE}های بزرگ. اگر جهت تخت مورب باشد، برآوردها
به‌شدت همبسته‌اند \lr{(trade-off).}

\textbf{3.} \lr{SE} روش \lr{Wald} را برای یک نسبت دوجمله‌ای از مشتق دوم
\lr{NLL} به دست آورید. \textbf{پاسخ}
\(\mathrm{NLL}''(\theta) = y/\theta^2 + (n-y)/(1-\theta)^2\). در
\(\hat\theta = y/n\) این برابر
\(n/\hat\theta + n/(1-\hat\theta) = n/[\hat\theta(1-\hat\theta)]\) است.
معکوس کردن و جذر گرفتن می‌دهد
\(\mathrm{SE} = \sqrt{\hat\theta(1-\hat\theta)/n}\).

\textbf{4.} مراحل یک \lr{bootstrap} پارامتری برای بازه‌ی اطمینان یک
پارامتر مدل را بنویسید. \textbf{پاسخ} مدل را برازش کنید تا
\(\hat{\boldsymbol\theta}\) به دست آید. \(B\) مجموعه‌داده از مدل با
\(\hat{\boldsymbol\theta}\) و طرح آزمایش اصلی شبیه‌سازی کنید. هرکدام را
دوباره برازش کنید. انحراف معیار برآوردهای جدید را \lr{SE} و صدک‌های 2.5 و
97.5 آن‌ها را بازه‌ی 95\% در نظر بگیرید.

\textbf{5.} چرا بازه‌ی \lr{profile likelihood} می‌تواند نامتقارن باشد ولی
بازه‌ی \lr{Wald} نمی‌تواند؟ \textbf{پاسخ} بازه‌ی \lr{Wald} یک \lr{NLL}
مربعی (توزیع نمونه‌گیری نرمال) فرض می‌کند که ذاتاً متقارن است. بازه‌ی
\lr{profile} از درست‌نمایی واقعی استفاده می‌کند که ممکن است چوله باشد،
مثلاً نزدیک یک کران یا برای پارامترهای نرخ.

\textbf{6.} بیست شرکت‌کننده هرکدام با نرخ متفاوت به‌طور نمایی یاد می‌گیرند.
منحنی میانگین با تابع توانی بهتر برازش می‌شود. آیا باید نتیجه بگیریم که
یادگیری از قانون توانی پیروی می‌کند؟ \textbf{پاسخ} نه. شکل توانی خطای
ناشی از میانگین‌گیری است: میانگین توابع نمایی با نرخ‌های ناهمگن منحنی‌ای
می‌سازد که نرخ نسبی تغییرش با زمان کم می‌شود. نتیجه‌گیری درباره‌ی شکل
یادگیری فردی باید از برازش‌های فردی (یا سلسله‌مراتبی) بیاید.

\textbf{7.} یک مزیت و یک عیب برازش جداگانه‌ی هر شرکت‌کننده بگویید.
\textbf{پاسخ} مزیت: تفاوت‌های فردی را رعایت می‌کند و از خطای میانگین‌گیری
دوری می‌کند. عیب: با \lr{trial}های کم برای هر نفر، برآوردها پرنویزند (و
گاهی حدی یا روی کران)، و تحلیل آن‌ها در مرحله‌ی دوم عدم قطعیت‌های متفاوتشان
را نادیده می‌گیرد.

\begin{center}\rule{0.5\linewidth}{0.5pt}\end{center}

\subsection{تکلیف}

این هفته تکلیف نمره‌دار جدیدی نداریم. \textbf{\lr{A2 (LS} در برابر
\lr{MLE)} در ابتدای درس هفته‌ی 6 تحویل داده می‌شود.} بخش \lr{A lab} امروز
نشان می‌دهد چطور به برآوردهای \lr{A2} خود خطای استاندارد اضافه کنید. این
کار الزامی نیست ولی تفسیرتان را قوی‌تر می‌کند.

\textbf{تمرین اختیاری (بدون نمره):} تمرین 5 (پوشش) بهترین آمادگی برای
مطالعات بازیابی هفته‌ی 7 است.

\begin{center}\rule{0.5\linewidth}{0.5pt}\end{center}

\subsection{مطالعه‌ی بیشتر}

\begin{itemize}
\tightlist
\item
  \textbf{\lr{Farrell} \lr{\&} \lr{Lewandowsky} (2018)، فصل 5.} ترکیب اطلاعات
  چند شرکت‌کننده: برازش تجمیعی در برابر فردی و خطرات میانگین‌گیری.
  \emph{(منبع اصلی.)}
\item
  \textbf{\lr{Farrell} \lr{\&} \lr{Lewandowsky} (2018)، فصل 4} (بخش‌های مربوط
  به عدم قطعیت پارامتر و \lr{Hessian)}، مرور از هفته‌ی قبل.
\item
  \emph{اختیاری، خارج از فهرست درس:} \lr{Heathcote, A., Brown, S.,} \lr{\&}
  \lr{Mewhort, D. J. K.} (2000). \lr{The power law repealed: The case
  for an exponential law of practice.} \emph{\lr{Psychonomic Bulletin}
  \lr{\&} \lr{Review,} 7}, 185\lr{--}207. تحلیل کلاسیک خطاهای میانگین‌گیری در
  داده‌های تمرین.
\end{itemize}

\emph{پیش‌نمایش هفته‌ی 6:} تمرین 4 هفته‌ی 3 نشان داد که یک مدل سه‌پارامتری
در \lr{SSE} از مدل واقعی جلو می‌زند. هفته‌ی بعد مستقیماً با این روبه‌رو
می‌شویم: چرا برازش خوب کافی نیست، و چطور \lr{cross-validation} چیزی را که
واقعاً می‌خواهیم، یعنی پیش‌بینی داده‌ی جدید، اندازه می‌گیرد.
